% ============================================================
\chapter{Probabilit\'e Conditionnelle et Ind\'ependance}
\label{ch:conditionnelle}
% ============================================================

Un m\'edecin apprend qu'un patient a un test positif pour une maladie rare.
Quelle est la probabilit\'e que le patient soit r\'eellement malade~? La
r\'eponse --- souvent contre-intuitive --- d\'epend crucialement de la
\emph{pr\'evalence} de la maladie et de la qualit\'e du test. C'est le
c\'el\`ebre \og paradoxe\fg{} de Bayes, et pour le r\'esoudre, il faut
ma\^itriser la notion de \emph{probabilit\'e conditionnelle}~: comment
mettre \`a jour nos croyances lorsqu'une information nouvelle est
disponible.

\begin{intuition}
La notion de probabilit\'e conditionnelle permet de mettre \`a jour l'information
disponible lorsqu'un \'ev\'enement est observ\'e. L'ind\'ependance, quant \`a elle,
formalise l'id\'ee que deux \'ev\'enements n'ont aucune influence l'un sur l'autre.
Ces deux concepts sont au c\oe ur de toute la th\'eorie des probabilit\'es.
\end{intuition}

% ------------------------------------------------------------
\section{Probabilit\'e conditionnelle}
% ------------------------------------------------------------

\begin{definition}[Probabilit\'e conditionnelle]
Soit $(\Omega, \mathcal{F}, \PP)$ un espace probabilis\'e et $B \in \mathcal{F}$
un \'ev\'enement tel que $\PP(B) > 0$. La \textbf{probabilit\'e conditionnelle}
de $A$ sachant $B$ est~:
\[
    \PP(A \mid B) = \frac{\PP(A \cap B)}{\PP(B)}.
\]
\end{definition}

\begin{proposition}
Pour $B$ fix\'e avec $\PP(B) > 0$, l'application $A \mapsto \PP(A \mid B)$
est une mesure de probabilit\'e sur $(\Omega, \mathcal{F})$.
\end{proposition}

\begin{proof}
V\'erifions les axiomes de Kolmogorov~:
\begin{itemize}
    \item $\PP(A \mid B) = \PP(A \cap B)/\PP(B) \geq 0$ pour tout $A$.
    \item $\PP(\Omega \mid B) = \PP(B)/\PP(B) = 1$.
    \item Si $(A_n)$ sont mutuellement disjoints, alors les $(A_n \cap B)$ le sont
    aussi, et par $\sigma$-additivit\'e de $\PP$~:
    \[
        \PP\!\left(\bigcup_n A_n \;\middle|\; B\right)
        = \frac{\PP\bigl(\bigcup_n (A_n \cap B)\bigr)}{\PP(B)}
        = \frac{\sum_n \PP(A_n \cap B)}{\PP(B)}
        = \sum_n \PP(A_n \mid B). \qedhere
    \]
\end{itemize}
\end{proof}

\begin{exemple}
On lance deux d\'es \'equilibr\'es. Sachant que la somme est 8, quelle est la
probabilit\'e que le premier d\'e affiche 3~?

Soit $A = \{\text{premier d\'e} = 3\}$ et $B = \{\text{somme} = 8\}$.
Les r\'esultats donnant somme 8 sont~: $(2,6), (3,5), (4,4), (5,3), (6,2)$,
donc $\PP(B) = 5/36$. Comme $A \cap B = \{(3,5)\}$, on a $\PP(A \cap B) = 1/36$.
Donc $\PP(A \mid B) = (1/36)/(5/36) = 1/5$.
\end{exemple}

% ------------------------------------------------------------
\section{R\`egle de la cha\^ine et formule de Bayes}
% ------------------------------------------------------------

\begin{theoreme}[R\`egle de multiplication / r\`egle de la cha\^ine]
Pour tous \'ev\'enements $A_1, \ldots, A_n$ avec
$\PP(A_1 \cap \cdots \cap A_{n-1}) > 0$~:
\[
    \PP(A_1 \cap \cdots \cap A_n)
    = \PP(A_1)\, \PP(A_2 \mid A_1)\, \PP(A_3 \mid A_1 \cap A_2)
    \cdots \PP(A_n \mid A_1 \cap \cdots \cap A_{n-1}).
\]
\end{theoreme}

\begin{definition}[Syst\`eme complet d'\'ev\'enements]
Une famille $(B_i)_{i \in I}$ d'\'ev\'enements forme un \textbf{syst\`eme complet
d'\'ev\'enements} (ou partition de $\Omega$) si~:
\begin{enumerate}[label=(\roman*)]
    \item $B_i \cap B_j = \emptyset$ pour $i \neq j$ (disjoints)\,;
    \item $\bigcup_{i \in I} B_i = \Omega$ (exhaustifs)\,;
    \item $\PP(B_i) > 0$ pour tout $i$.
\end{enumerate}
\end{definition}

\begin{theoreme}[Formule des probabilit\'es totales]
\label{thm:prob_totales}
Soit $(B_i)_{i \in I}$ un syst\`eme complet d'\'ev\'enements (avec $I$ fini ou
d\'enombrable). Pour tout $A \in \mathcal{F}$~:
\[
    \PP(A) = \sum_{i \in I} \PP(A \mid B_i)\, \PP(B_i).
\]
\end{theoreme}

\begin{proof}
Comme $A = A \cap \Omega = \bigcup_{i} (A \cap B_i)$ (union disjointe)~:
\[
    \PP(A) = \sum_{i} \PP(A \cap B_i) = \sum_{i} \PP(A \mid B_i)\, \PP(B_i). \qedhere
\]
\end{proof}

\begin{theoreme}[Formule de Bayes]
\label{thm:bayes}
Soit $(B_i)_{i \in I}$ un syst\`eme complet d'\'ev\'enements et $A$ un \'ev\'enement
avec $\PP(A) > 0$. Alors pour tout $j \in I$~:
\[
    \PP(B_j \mid A) = \frac{\PP(A \mid B_j)\, \PP(B_j)}{\sum_{i \in I} \PP(A \mid B_i)\, \PP(B_i)}.
\]
\end{theoreme}

\begin{intuition}[title=\textbf{Interpr\'etation bay\'esienne}]
Les $\PP(B_i)$ sont les probabilit\'es \emph{a priori} des hypoth\`eses $B_i$.
Apr\`es observation de $A$, la formule de Bayes fournit les probabilit\'es
\emph{a posteriori} $\PP(B_i \mid A)$, mises \`a jour par la vraisemblance
$\PP(A \mid B_i)$.
\end{intuition}

\begin{exemple}[Test de d\'epistage]
Un test m\'edical a une sensibilit\'e de 99\% (vrai positif) et une sp\'ecificit\'e
de 95\% (vrai n\'egatif). La pr\'evalence de la maladie est de 0{,}1\%. Quelle est la
probabilit\'e qu'un patient soit malade sachant que le test est positif~?

Soit $M$ = \og malade \fg et $T^+$ = \og test positif \fg. On a~:
$\PP(M) = 0{,}001$, $\PP(T^+ \mid M) = 0{,}99$, $\PP(T^+ \mid M^c) = 0{,}05$.
\begin{align*}
    \PP(M \mid T^+)
    &= \frac{\PP(T^+ \mid M)\, \PP(M)}{\PP(T^+ \mid M)\, \PP(M) + \PP(T^+ \mid M^c)\, \PP(M^c)} \\
    &= \frac{0{,}99 \times 0{,}001}{0{,}99 \times 0{,}001 + 0{,}05 \times 0{,}999}
    = \frac{0{,}00099}{0{,}05094}
    \approx 0{,}0194.
\end{align*}
M\^eme avec un test tr\`es performant, la probabilit\'e d'\^etre malade sachant un
test positif n'est que d'environ 2\%, \`a cause de la faible pr\'evalence.
\end{exemple}

% ------------------------------------------------------------
\section{Ind\'ependance d'\'ev\'enements}
% ------------------------------------------------------------

\begin{definition}[Ind\'ependance de deux \'ev\'enements]
Deux \'ev\'enements $A$ et $B$ sont \textbf{ind\'ependants}, not\'e
$A \indep B$, si~:
\[
    \PP(A \cap B) = \PP(A) \cdot \PP(B).
\]
\end{definition}

\begin{remarque}
Si $\PP(B) > 0$, l'ind\'ependance \'equivaut \`a $\PP(A \mid B) = \PP(A)$~:
savoir que $B$ s'est r\'ealis\'e ne modifie pas la probabilit\'e de $A$.
\end{remarque}

\begin{attention}[title=\textbf{Ind\'ependance $\neq$ incompatibilit\'e}]
Deux \'ev\'enements disjoints ($A \cap B = \emptyset$) de probabilit\'e non nulle
ne sont \emph{jamais} ind\'ependants, car $\PP(A \cap B) = 0 \neq \PP(A)\PP(B)$.
Ne confondez pas ces deux notions~!
\end{attention}

\begin{definition}[Ind\'ependance mutuelle]
Les \'ev\'enements $A_1, \ldots, A_n$ sont \textbf{mutuellement ind\'ependants}
si, pour tout sous-ensemble $J \subseteq \{1, \ldots, n\}$ avec $\abs{J} \geq 2$~:
\[
    \PP\!\left(\bigcap_{j \in J} A_j\right) = \prod_{j \in J} \PP(A_j).
\]
\end{definition}

\begin{attention}[title=\textbf{Ind\'ependance deux \`a deux $\neq$ mutuelle}]
L'ind\'ependance deux \`a deux ne suffit pas pour l'ind\'ependance mutuelle.
Il faut v\'erifier \emph{toutes} les sous-familles, pas seulement les paires.
\end{attention}

\begin{exemple}[Contre-exemple de Bernstein]
On lance deux d\'es \'equilibr\'es. Posons~:
\begin{itemize}
    \item $A$ = \og le premier d\'e est pair \fg,
    \item $B$ = \og le second d\'e est pair \fg,
    \item $C$ = \og la somme est paire \fg.
\end{itemize}
On v\'erifie que $\PP(A) = \PP(B) = \PP(C) = 1/2$ et que toute paire est
ind\'ependante~: $\PP(A \cap B) = 1/4 = \PP(A)\PP(B)$, etc. Cependant,
$\PP(A \cap B \cap C) = 1/4 \neq 1/8 = \PP(A)\PP(B)\PP(C)$.
Les trois \'ev\'enements sont ind\'ependants deux \`a deux mais pas mutuellement.
\end{exemple}

\begin{proposition}
Si $A$ et $B$ sont ind\'ependants, alors les couples suivants le sont aussi~:
$(A, B^c)$, $(A^c, B)$, $(A^c, B^c)$.
\end{proposition}

\begin{proof}
$\PP(A \cap B^c) = \PP(A) - \PP(A \cap B) = \PP(A) - \PP(A)\PP(B)
= \PP(A)(1 - \PP(B)) = \PP(A)\PP(B^c)$.
Les autres cas sont analogues.
\end{proof}

% ------------------------------------------------------------
\section{Second lemme de Borel--Cantelli}
% ------------------------------------------------------------

\begin{theoreme}[Second lemme de Borel--Cantelli]
\label{thm:borel_cantelli_2}
Soit $(A_n)_{n \geq 1}$ une suite d'\'ev\'enements \textbf{mutuellement
ind\'ependants}. Si $\sum_{n=1}^{\infty} \PP(A_n) = +\infty$, alors~:
\[
    \PP\!\left(\limsup_{n \to \infty} A_n\right) = 1.
\]
Autrement dit, presque s\^urement, une infinit\'e de $A_n$ se r\'ealisent.
\end{theoreme}

\begin{proof}
Il suffit de montrer que $\PP\bigl(\bigcup_{k=n}^{\infty} A_k\bigr) = 1$
pour tout $n$. Passant au compl\'ementaire~:
\[
    \PP\!\left(\bigcap_{k=n}^{N} A_k^c\right)
    = \prod_{k=n}^{N} \PP(A_k^c)
    = \prod_{k=n}^{N} (1 - \PP(A_k))
    \leq \prod_{k=n}^{N} e^{-\PP(A_k)}
    = \exp\!\left(-\sum_{k=n}^{N} \PP(A_k)\right).
\]
Quand $N \to \infty$, la somme diverge, donc le produit tend vers 0. Par
continuit\'e d\'ecroissante~:
$\PP\bigl(\bigcap_{k=n}^{\infty} A_k^c\bigr) = 0$, d'o\`u
$\PP\bigl(\bigcup_{k=n}^{\infty} A_k\bigr) = 1$.
\end{proof}

\begin{formules}[title=\textbf{R\'esum\'e : Borel--Cantelli}]
\begin{itemize}
    \item $\sum \PP(A_n) < \infty \implies \PP(\limsup A_n) = 0$ \quad (sans hypoth\`ese).
    \item $\sum \PP(A_n) = \infty$ \textbf{et} ind\'ependance
          $\implies \PP(\limsup A_n) = 1$.
\end{itemize}
\end{formules}

% ------------------------------------------------------------
\section{Ind\'ependance de tribus et de variables al\'eatoires}
% ------------------------------------------------------------

\begin{definition}[Tribus ind\'ependantes]
Deux sous-tribus $\mathcal{G}_1, \mathcal{G}_2 \subseteq \mathcal{F}$ sont
\textbf{ind\'ependantes} si, pour tous $A \in \mathcal{G}_1$ et
$B \in \mathcal{G}_2$~:
\[
    \PP(A \cap B) = \PP(A)\, \PP(B).
\]
\end{definition}

\begin{definition}[Variables al\'eatoires ind\'ependantes --- pr\'eambule]
Deux variables al\'eatoires $X$ et $Y$ sont \textbf{ind\'ependantes} si les
tribus qu'elles engendrent, $\sigma(X)$ et $\sigma(Y)$, sont ind\'ependantes.
En termes de lois (voir chapitre~3 pour les d\'efinitions compl\`etes),
cela revient \`a~:
\[
    \PP(X \in A, Y \in B) = \PP(X \in A)\, \PP(Y \in B)
\]
pour tous bor\'eliens $A, B \in \mathcal{B}(\R)$.
\end{definition}

% ------------------------------------------------------------
\section{Applications classiques}
% ------------------------------------------------------------

\begin{exemple}[Probl\`eme de Monty Hall]
Un candidat choisit une porte parmi trois. Derri\`ere l'une se trouve une voiture,
derri\`ere les deux autres des ch\`evres. Le pr\'esentateur, qui sait o\`u est la
voiture, ouvre une porte r\'ev\'elant une ch\`evre, puis propose au candidat de
changer de porte.

Soit $V_i$ = \og la voiture est derri\`ere la porte $i$ \fg ($i=1,2,3$), et
supposons que le candidat choisisse la porte~1. Par la formule de Bayes, on
montre que la probabilit\'e de gagner en changeant est $2/3$, tandis que
rester donne $1/3$. Il est donc avantageux de changer.
\end{exemple}

\begin{exemple}[Marche al\'eatoire et ruine du joueur]
Un joueur commence avec un capital $k \in \{1, \ldots, N-1\}$. \`A chaque tour,
il gagne 1\,\texteuro{} avec probabilit\'e $p$ ou perd 1\,\texteuro{} avec probabilit\'e $q = 1-p$,
ind\'ependamment. Le jeu s'arr\^ete quand le capital atteint 0 (ruine) ou $N$
(objectif atteint).

En utilisant la probabilit\'e conditionnelle et un argument de r\'ecurrence, on
montre que la probabilit\'e de ruine est~:
\[
    \PP(\text{ruine} \mid X_0 = k) =
    \begin{cases}
        \dfrac{(q/p)^k - (q/p)^N}{1 - (q/p)^N} & \text{si } p \neq q, \\[6pt]
        1 - \dfrac{k}{N} & \text{si } p = q = \tfrac{1}{2}.
    \end{cases}
\]
\end{exemple}

% ------------------------------------------------------------
\section{Exercices}
% ------------------------------------------------------------

\begin{exercice}
Un lot contient 10 pi\`eces dont 3 sont d\'efectueuses. On tire 2 pi\`eces
successivement sans remise. Calculer la probabilit\'e que la seconde soit
d\'efectueuse sachant que la premi\`ere l'est.
\end{exercice}

\begin{exercice}
Trois machines $M_1, M_2, M_3$ produisent respectivement 50\%, 30\% et 20\%
de la production. Leurs taux de d\'efaut sont 2\%, 3\% et 5\%. Un article pris au
hasard est d\'efectueux. Quelle est la probabilit\'e qu'il provienne de $M_2$~?
\end{exercice}

\begin{exercice}
Montrer que si $A$, $B$, $C$ sont mutuellement ind\'ependants, alors $A \cup B$
et $C$ sont ind\'ependants.
\end{exercice}

\begin{exercice}
On lance une pi\`ece \'equilibr\'ee $n$ fois. Soit $A_k$ l'\'ev\'enement
\og le $k$-i\`eme lancer donne Pile \fg. Montrer que les $A_k$ sont mutuellement
ind\'ependants.
\end{exercice}

\begin{exercice}[Formule de Bayes it\'er\'ee]
On dispose d'une urne contenant 1 boule rouge et 1 boule bleue. On tire une boule
au hasard, on la remet et on ajoute une boule de la m\^eme couleur. Apr\`es $n$
tirages, calculer la probabilit\'e que toutes les boules tir\'ees soient rouges.
\end{exercice}

\begin{exercice}
Soit $(A_n)_{n \geq 1}$ une suite d'\'ev\'enements ind\'ependants avec
$\PP(A_n) = 1/(n+1)$. Montrer que $\PP(\limsup A_n) = 1$. Interpr\'eter.
\end{exercice}

\begin{exercice}[Paradoxe de Simpson]
Donner un exemple num\'erique o\`u $\PP(A \mid B) > \PP(A \mid B^c)$
mais $\PP(A \mid B \cap C) < \PP(A \mid B^c \cap C)$ et
$\PP(A \mid B \cap C^c) < \PP(A \mid B^c \cap C^c)$. Expliquer le ph\'enom\`ene.
\end{exercice}

\begin{exercice}
Soit $p \in (0,1)$ et $(X_n)_{n \geq 1}$ une suite de lancers de Bernoulli
ind\'ependants de param\`etre $p$. Montrer que $\PP(\exists n : X_n = 1) = 1$.
\end{exercice}
